\subsection{All-height split class groups in the two-prime tower}

\begin{statement}{A.177}{Proved}
{Every three-class group is generated by the six changed primes}
Retain $d=649$, $K_n=\Q(\rho_n)$, $\rho_n^{3^n}=3d^{24}$,
$E_n=K_n(\zeta_3)$ and $S$ above $2,3,11,59,\infty$.
For every $n\ge2$, the following assertions hold:
\[
 \begin{gathered}
 \operatorname{Cl}(E_n)\{3\}
   =\langle[R]:R\mid11\text{ or }59\rangle,\\
 \operatorname{Cl}_S(E_n)\{3\}=0,
       \qquad \operatorname{Cl}_S(K_n)\{3\}=0,\\
 [\mathcal O_{E_n,S}^{\times}:
   \mathcal O_{E_n,S}^{\times}\cap N_{E_{n+1}/E_n}E_{n+1}^{\times}]
                         =3^{2n+6}.
 \end{gathered}
\tag{A.177.1}
\]
There are exactly six primes in the first generating set.
The actual action of $\operatorname{Gal}(E_{n+1}/E_n)$ on
$\operatorname{Cl}(E_{n+1})\{3\}$ is trivial.
The proof constructs auxiliary norm targets with prescribed two-prime
divisors; it does not assert that those targets preserve an arbitrarily
prescribed earlier norm representative.
\end{statement}

\begin{proof}
\emph{The local factors at every height.}
Each step $E_{n+1}/E_n$ is cyclic cubic, with
$\rho_{n+1}^3=\rho_n$ and actual norm $N\rho_{n+1}=\rho_n$.
At each of the three primes above $q=11,59$ the residue field is
$\F_{q^2}$, the valuation of $\rho_n$ is eight, and the next step
is totally ramified of degree three.  At the unique three-prime
its valuation is two and the next step is again totally ramified.
The root equations give no other ramification.

At two, the factors of $K_n$ have residue degrees
\[
                  1,\ 2,\ 2\cdot3,\ldots,2\cdot3^{n-1}.
\tag{A.177.2}
\]
These are the actual unramified factors of successive roots of a
two-adic unit.  Its unique roots in the principal-unit group separate
off the roots of unity; the nontrivial residues have orders
$3,\ldots,3^n$, and $\operatorname{ord}_{3^j}(2)=2\cdot3^{j-1}$.
For the latter identity, $v_3(4-1)=1$ and the binomial expansion
of $(1+3^{r+1}u)^3-1$, with $3\nmid u$, has unique lowest
valuation $r+2$.  Induction gives
$v_3(4^{3^r}-1)=r+1$; together with the order two of $2$ modulo
three, this proves the stated order at every height.
Adjoining $\zeta_3$ gives one degree-two completion above the first
factor and two completions above each of the remaining factors.
Thus $E_n$ has $2n+1$ two-primes.  Its distinguished first prime
splits into three primes in $E_{n+1}$, and each of the other $2n$
primes is unramified inert of degree three.  In particular, the
number of nonsplit places in $S$ is exactly
\[
                          2n+6+1=2n+7.
\tag{A.177.3}
\]
All infinite places are complex.  The full principal divisor of the
chosen radical is
\[
 \operatorname{div}_{E_n}(\rho_n)
      =2\mathfrak p_{3,n}+8\sum_{R\mid11,59}R .
\tag{A.177.4}
\]
This follows directly from the ramification indices
$2\cdot3^n$ at three and $3^{n-1}$ at the changed primes.

\emph{A fixed six-dimensional local norm image.}
Use the nine literal global units of A.175.8 and adjoin
\[
                         h_{11}=\frac{11}{(2+\alpha)^3}.
\tag{A.177.5}
\]
Every one of these ten fixed elements of $E_2$ has valuation zero
at every two-prime and at three.  At level two, the six-row matrix
of the nine units is A.175.10.  The normalized tame norm column
of $h_{11}$ in that row order is
\[
                             (0,1,2,0,0,0)^t.
\tag{A.177.6}
\]
For the two quadratic eleven-branches, put $\varpi=t^2/11$.
Then $\varpi^3=11(d/11)^4\alpha^2$, and $h_{11}$ modulo a
cube has normalized unit $\alpha^{-2}$, up to rational units.
Rational residue units are cubes in $\F_{11^2}$; its logarithms
are one and two.  At the linear branch $\alpha$ is rational in
the residue field and the logarithm is zero.  At fifty-nine the
numerator is rational and the denominator is already a cube, so
all three logs are zero.  This proves A.177.6 with the actual
uniformizer and residue maps.  The ten-column matrix has rank six.

This same rank persists at every level.  Let $x$ be any of the
ten fixed elements and follow one of the unique changed-prime
chains.  The cubic norm character for the next step is the local
Hilbert pairing with $\rho_n$.  The projection formula and the
actual equality $N\rho_{n+1}=\rho_n$ give
\[
 \operatorname{inv}_{E_{n+1,R}}
       (\kappa_3(x)\cup\kappa_3(\rho_{n+1}))
 =\operatorname{inv}_{E_{n,R}}
       (\kappa_3(x)\cup\kappa_3(\rho_n)).
\tag{A.177.7}
\]
Here corestriction is applied to the left-hand cup product;
local invariant is unchanged under corestriction.  The common
$\mu_3$ and its chosen generator identify the coefficient twists.
The equality is the literal norm projection formula, and does not
identify characters merely from their orders.
It proves that the ten fixed $S$-units give a surjective map to
the six local norm quotients at eleven and fifty-nine at every
height, while their two-prime valuations remain zero.

\emph{The induction step for the $S$-class group.}
A.173 supplies the initial assertion at $n=2$: its whole ordinary
three-class group is generated by changed-prime classes.
Suppose this assertion has been proved at height $n$, and write
$A_n=\operatorname{Cl}(E_n)\{3\}$.
For each of the $2n$ growing two-primes $P$, choose an integral
linear combination $D_P$ of the six changed primes with
$[P-D_P]_{(3)}=0$.
The class of that fractional divisor has finite order $h_P$ prime
to three.  There is therefore an actual element $x_P$ such that
\[
                    \operatorname{div}(x_P)=h_P(P-D_P).
\tag{A.177.8}
\]
It is an $S$-unit.  Its valuations at the growing two-primes are
zero except for $h_P$ at $P$.
Since those local extensions are unramified cubic, these $2n$
valuation rows form a diagonal invertible block modulo three.
The ten elements of A.177.5--7 supply an independent six-row
changed-prime block with all two-prime valuations zero.
The resulting local norm image has rank $2n+6$.

All other local norm characters of an $S$-unit vanish outside
$S$.  There are $2n+7$ nonsplit places by A.177.3, each with
local quotient $\Z/3$.  Global reciprocity places the image in
the kernel of their sum map, whose dimension is $2n+6$.
The cyclic Hasse norm theorem identifies the kernel of these
local characters with actual field norms
\cite[Chapter VIII, Theorem 3.1, printed p.~234]{MilneCFT}.
The lower and upper bounds agree and prove the index in A.177.1.

The induction hypothesis also gives
$\operatorname{Cl}_S(E_n)\{3\}=0$.
Apply the fully derived $S$-ambiguous formula A.175.18: all
ramification lies inside $S$, and its decomposition-degree product
is $3^{2n+7}$.  It gives
\[
 \left|\bigl(\operatorname{Cl}_S(E_{n+1})\{3\}\bigr)
                    ^{\operatorname{Gal}(E_{n+1}/E_n)}\right|
               =\frac{3^{2n+7}}{3\,3^{2n+6}}=1.
\tag{A.177.9}
\]
The finite three-group fixed-point argument of A.175 proves
$\operatorname{Cl}_S(E_{n+1})\{3\}=0$.

\emph{An auxiliary actual norm that controls the split two-prime orbit.}
Let $P_0$ be the distinguished two-prime of $E_n$.
Choose a changed-prime divisor $D$ and $h$ prime to three with
\[
                         \operatorname{div}(x)=h(P_0-D).
\tag{A.177.10}
\]
The same principal-power argument as in A.177.8 gives this actual
element.  Multiply $x$ by a product of the ten fixed elements
of A.177.5--7, with integer exponents, so that all its six
changed-prime norm characters vanish.  Call the result $x'$.
Its two-divisor is still exactly $hP_0$ and its three-valuation
is still zero.
At $P_0$ the next local algebra is split and its norm is onto.
At every growing two-prime, $x'$ is a unit in an unramified
cubic extension and is a local norm.  It is a local norm at
the changed primes by construction and outside $S$ by the unit
criterion.  Reciprocity makes its final character at three zero.
The cyclic Hasse norm theorem gives a field-norm preimage.

We retain the actual support repair.  If $z\in E_{n+1}^\times$
has norm $x'$, its outside-$S$ divisor has norm zero and equals
$(\sigma-1)J$ on the unramified orbit divisor lattice.
Choose a prime-to-three integer $a$ killing $[J]$ in the
$S$-class group, which has zero three-primary part by A.177.9,
and choose $v$ with $\operatorname{div}_{S^c}(v)=aJ$.
Then $z^a/(\sigma v/v)$ is an actual $S$-unit of norm $(x')^a$.
Choose integers $b,c$ with $ab+3c=1$ and multiply its $b$-th
power by $(x')^c$.  This produces an actual $S$-unit $y$ with
\[
                             N_{E_{n+1}/E_n}y=x'.
\tag{A.177.11}
\]
This is precisely the finite divisor and scalar correction of
A.176.7--9, applied to the present auxiliary target.
No compatibility with a separately prescribed lower norm is used.

\emph{The polynomial that forces ordinary class generation.}
Let $B_{n+1}$ be the subgroup of $A_{n+1}$ generated by its six
changed-prime classes, and put $C=A_{n+1}/B_{n+1}$.
The proved $S$-class vanishing makes $A_{n+1}$ generated by the
classes of its $S$-primes.
Equation A.177.4 kills the three-prime in $C$, because its
coefficient two is invertible on a three-group.
Every growing upper two-prime is the ideal extension of its
unique lower prime, with ramification weight one.  The lower
prime class is a combination of changed-prime classes, whose
extensions are three times their unique upper changed primes.
These growing two-prime classes consequently vanish in $C$.
Thus $C$ is generated by the three primes above $P_0$.

Write them as $P_i=\sigma^iP_0'$ for $i=0,1,2$.
The norm in A.177.11 and its exact two-divisor give integers
$k_0,k_1,k_2$ with
\[
 \operatorname{div}(y)=\sum_{i=0}^2k_iP_i+D_q,
       \qquad k_0+k_1+k_2=h,
\tag{A.177.12}
\]
where $D_q$ is a changed-prime divisor.  Indeed every growing
two-prime has norm coefficient three and target valuation zero,
and the unique three-prime has norm residue weight one and target
valuation zero.  There is no other support.
Consequently $C$ is a cyclic $\Z_3[\langle\sigma\rangle]$-module
generated by $c_0=[P_0']$, and
\[
          (k_0+k_1\sigma+k_2\sigma^2)c_0=0,
                    \qquad k_0+k_1+k_2\not\equiv0\pmod3.
\tag{A.177.13}
\]
On $C/3C$, the operator $\sigma-1$ has cube zero.
The polynomial in A.177.13 is a nonzero scalar plus a nilpotent
polynomial in that operator, and hence is invertible.
It kills a module generator, so $C/3C=0$.
A finite three-group equal to three times itself is zero.
Thus $C=0$, completing the induction of ordinary class generation.

Every generating changed prime is totally ramified in the next
cyclic step, so its unique upper prime is fixed by that step's
Galois group.  Applied to the generation just proved at height
$n+1$, this proves triviality of the stated action on $A_{n+1}$.
Finally quadratic ideal extension and norm between $K_n$ and
$E_n$ have composite multiplication by two, also on $S$-class
groups.  Extension is injective on the three-primary part.
The vanishing for $E_n$ therefore proves the asserted vanishing
for $K_n$.
\end{proof}

The exact fixed local-character rank, the residue-degree counts,
the auxiliary divisor norm and the finite polynomial-invertibility
steps are checked in
\path{code/two_inert_all_height_split_certificate.py}, with retained
output and data.  The proof is an induction through actual finite
fields and finite class groups; the certificate does not claim to
execute infinitely many class-group calculations.
